% Chapter 1 Introduction

\chapter[Introduction]{\centering Introduction}\label{chap1}
\newpage

\section{Background}
IT procurement is an important part of an organization’s operations as it involves the acquisition and management of hardware, software licenses, SaaS subscriptions, and related services. The procurement process involves several stages including vendor identification, document processing, approval, inventory management, contract handling and vendor risk assessment. In many organizations, these activities are handled through emails, spreadsheets, invoices, and separate systems, making it challenging to maintain a consistent and updated view of procurement activities. As a result, approvals can be delayed, purchases duplicated and licenses left unused, difficulties with tracking contracts and increased vendor-related risks. The opportunity to use artificial intelligence and automation presents an opportunity to improve these activities and make the procurement process more efficient.


\section{Motivation}
The rationale behind the project revolves around the need to reduce the amount of time spent on IT procurement manual processes as well as increase transparency and control during the procedure. Currently, several steps have to be taken by employees in the process, including verifying documents, reviewing vendor details, requesting and obtaining approvals, checking the utilisation of inventory, and ensuring the renewal of contracts among other aspects. These procedures create difficulties, especially when the data involved is high in volume, and a system that can help automate and optimise processing and support cross-functional communication would be of great benefit. The project, therefore, aims to implement a set of special-purpose software agents that will utilise artificial intelligence to assist with document processing, vendor identification, approvals, inventory use analysis, contract management, and risk detection among other fields.
\section{Aim and Objectives}

\subsection{Aim}

To develop an \textbf{Agentic AI-Based IT Procurement Intelligence Platform} that integrates, automates, and intelligently manages major activities involved in the IT procurement lifecycle while maintaining human oversight for critical decisions.

\subsection{Objectives}

\begin{itemize}
    \item To process and extract relevant information from procurement documents such as purchase orders, invoices, and quotations.
    
    \item To verify vendor information and identify duplicate vendors and invoices.
    
    \item To automate approval workflows based on spending limits, budgets, user roles, and inventory availability.
    
    \item To identify underutilised software licenses and support their controlled reclamation.
    
    \item To assist in contract management and monitor important renewal-related information.
    
    \item To continuously assess vendor-related risks using machine learning techniques.
    
    \item To provide shared platform services such as notifications, authentication, monitoring, and system management.
    
    \item To maintain human involvement in low-confidence results and critical procurement decisions.
\end{itemize}

\section{Problem statement}


IT procurement involves several activities such as collecting procurement documents, verifying vendor information, processing invoices and quotations, obtaining approvals, managing software licenses, monitoring contracts, and assessing vendor risks. These activities are often handled through multiple systems and manual processes, resulting in delays, duplicate records, inefficient resource utilisation, and limited visibility across the procurement lifecycle.

Traditional procurement systems primarily focus on recording and managing transactions rather than intelligently analysing procurement data and assisting users in decision-making. Identifying duplicate vendors or invoices, detecting underutilised software licenses, tracking contract renewals, and assessing vendor-related risks may require significant manual effort. Furthermore, automated decisions without proper confidence assessment and human oversight can lead to incorrect procurement actions.

Therefore, there is a need for an intelligent and integrated solution that can automate routine procurement activities, analyse relevant data, assist in approval and decision-making processes, and continuously monitor procurement-related risks. The proposed Agentic AI-Based IT Procurement Intelligence Platform aims to address these challenges by integrating AI-driven analysis, workflow automation, document processing, vendor intelligence, license management, contract monitoring, and human-in-the-loop mechanisms within a unified platform.

\section{Scope}

The scope of the proposed Agentic AI-Based IT Procurement Intelligence Platform includes the development of an integrated system to support major activities involved in the IT procurement lifecycle. The platform will focus on the following areas:

\begin{itemize}
    \item \textbf{Procurement Document Processing:} Processing purchase orders, invoices, quotations, and other relevant procurement documents to extract and structure important information.

    \item \textbf{Vendor Intelligence:} Verifying vendor information, identifying duplicate vendors, and maintaining relevant vendor records.

    \item \textbf{Invoice and Duplicate Detection:} Analysing invoice information to identify duplicate or potentially inconsistent records.

    \item \textbf{Approval Workflow Automation:} Supporting procurement approval workflows based on predefined spending limits, budgets, user roles, and inventory availability.

    \item \textbf{Software License Management:} Identifying underutilised software licenses and supporting their controlled reclamation or reassignment.

    \item \textbf{Contract Management:} Maintaining contract-related information and monitoring important dates, including contract renewal and expiration.

    \item \textbf{Vendor Risk Assessment:} Applying machine learning techniques to analyse vendor-related information and continuously assess potential risks.

    \item \textbf{Shared Platform Services:} Providing common services such as authentication, notifications, monitoring, logging, and system management.

    \item \textbf{Human-in-the-Loop Decision Making:} Routing low-confidence results and critical procurement decisions to authorised users for review and approval.

    \item \textbf{Agentic AI Integration:} Using AI agents to coordinate relevant procurement tasks, analyse information, and assist users across different stages of the procurement lifecycle.
\end{itemize}

The platform is intended to assist procurement teams by reducing manual effort, improving data accuracy, increasing process visibility, and supporting informed decision-making. The system will focus on decision support and workflow automation rather than completely replacing human decision-makers in critical procurement activities.
```

\section{Report Outline}



This report is organized into six chapters. Chapter 1, \textbf{Introduction}, presents the background and motivation of the project, followed by the aim and objectives, problem statement, scope, and report outline. Chapter 2, \textbf{Study of the System}, discusses the techniques relevant to the proposed system, various available techniques, and related works in the field of IT procurement and Agentic AI. Chapter 3, \textbf{Proposed System}, describes the problem addressed by the project, its scope, and the proposed Agentic AI-Based IT Procurement Intelligence Platform. Chapter 4, \textbf{Design of the System}, covers requirement engineering, requirement elicitation, software lifecycle model, requirement analysis, UML diagrams and DFDs, cost analysis, hardware and software requirements, system architecture, UI/UX design, and the block diagram of the system. Chapter 5, \textbf{Result and Discussion}, presents the implemented system through screenshots, important code and main logic, and discusses the testing performed, including unit testing, integration testing, and black-box testing. Chapter 6, \textbf{Conclusion and Future Scope}, summarizes the project outcomes, limitations, and possible future enhancements. The report also includes references and appendices covering the project timeline, SDG mapping, complex engineering problem identification and evaluation, and project outcomes.




%%%% End of Chapter 1